\section{Finite-strain homogenization and the reference problem}
\label{sec:fom}
\subsection{Kinematics and work-conjugate variables}
We consider an in-plane deformation gradient $\F\in\GL$, with
\begin{equation}
 \C=\F^T\F,\qquad J=\det\F>0,\qquad
 \E=\tfrac12(\C-\I).
\end{equation}
Plane strain means that the three-dimensional gradient is $\diag(\F,1)$; it does not mean that the out-of-plane stress vanishes. The engineering strain vector and its work-conjugate stress vector are
\begin{equation}
 \ee=(E_{11},E_{22},2E_{12})^T,\qquad
 \ssv=(S_{11},S_{22},S_{12})^T,\qquad
 \Sg:\dd\E=\ssv\cdot\dd\ee.
 \label{eq:voigt}
\end{equation}
We also write $\gamma_{12}=2E_{12}$. For an objective differentiable energy $W$, the second Piola stress and constitutive tangent are $\ssv=\partial_{\ee}W$ and $\bm D=\partial^2_{\ee\ee}W$. The factors in Eq.~\eqref{eq:voigt} are essential for both energy-cycle tests and derivative checks.

\subsection{Periodic microscopic equilibrium}
Let $Y$ denote a reference periodic cell and $Y_s\subset Y$ its solid part. A homogeneous macroscopic gradient $\overline{\F}$ induces the microscopic deformation
\begin{equation}
 \bm\varphi(\bm X)=\overline{\F}\bm X+\bm w(\bm X),\qquad
 \F_\mu=\overline{\F}+\nabla\bm w,
\end{equation}
where $\bm w$ is periodic on opposite faces. A constraint removes its translational null space, and the void surface is traction-free. On a differentiable equilibrium branch, the discrete fluctuation coordinates $\bm a^*(\overline{\F})$ satisfy
\begin{equation}
 \bm r(\bm a,\overline{\F})=\partial_{\bm a}\Pi_\mu(\bm a,\overline{\F})=\bm0,
 \qquad
 \Pi_\mu=\int_{Y_s}\psi_\mu(\F_\mu)\,\dd V.
 \label{eq:microeq}
\end{equation}
The branch energy and homogenized first Piola stress are
\begin{equation}
 \overline W=\frac{\Pi_\mu(\bm a^*,\overline{\F})}{|Y|},\qquad
 \overline{\bm P}=\frac{1}{|Y|}\int_{Y_s}\bm P_\mu\,\dd V,
 \qquad \overline{\Sg}=\overline{\F}^{-1}\overline{\bm P}.
 \label{eq:homstress}
\end{equation}
The denominator includes the void. Averaging microscopic second Piola stress instead of first Piola stress would generally give a different quantity and is not used. Stationarity in Eq.~\eqref{eq:microeq} eliminates the fluctuation derivative in the first derivative of $\overline W$. More explicitly, periodic test fluctuations satisfy
\begin{equation}
 \int_{Y_s}\bm P_\mu:\nabla\delta\bm w\,\dd V=0,
 \qquad
 \frac1{|Y|}\int_{Y_s}\bm P_\mu:\delta\F_\mu\,\dd V
     =\overline{\bm P}:\delta\overline{\F}.
 \label{eq:hill_mandel}
\end{equation}
The second identity follows from the first and the additive microkinematics. It is the work-consistency condition of the periodic micro-to-macro transition~\cite{Miehe1999}. Void boundaries contribute no traction work. The cell is an ideal periodic material model; no claim of statistical representativeness for a random microstructure is needed.

The implementation queries the constitutive law through $\overline{\F}=\sqrt{\I+2\E}$, exploiting objectivity to remove a superposed macroscopic rotation. The microscopic Newton method follows an equilibrium branch; convergence does not establish global minimization of $\Pi_\mu$. In particular, no general assertion that homogenization preserves polyconvexity is required here.

\subsection{Consistent tangent and the scope of the reference}
If $\bm r_{,\bm a}$ is nonsingular on the selected branch, implicit differentiation gives
\begin{equation}
 \frac{\dd\bm a^*}{\dd\ee}=-\bm r_{,\bm a}^{-1}\bm r_{,\ee},\qquad
 \bm D=\ssv_{,\ee}-\ssv_{,\bm a}\bm r_{,\bm a}^{-1}\bm r_{,\ee}.
 \label{eq:ift}
\end{equation}
Every derivative in this expression must refer to the stress and residual actually evaluated. The reference and reduced implementations are checked against perturbations that re-solve microscopic equilibrium. These checks validate linearization, not a change of homogenization assumptions.

The nested FOM--FE$^2$ solution is the reference for surrogate errors at fixed macro and micro discretizations. It is not a direct numerical simulation of all pores in the coupon. Homogenization error, boundary-layer effects, and mesh-discretization error therefore remain conceptually distinct from the surrogate errors reported below.

\subsection{Macroscopic equilibrium and the constitutive interface}
Let $\bm U$ collect the free macroscopic nodal displacements, and let $\ee_g(\bm U)$ be the engineering Green strain at macroscopic integration point $g$. With $\bm B_g=\partial\ee_g/\partial\bm U$, reference integration weights $\Omega_g$, and dead external loads $\bm f_{\rm ext}$, equilibrium reads
\begin{equation}
 \bm R_M(\bm U)=
       \sum_g\Omega_g\bm B_g^T\ssv_g(\ee_g(\bm U))
       -\bm f_{\rm ext}=\bm0.
 \label{eq:macro_residual}
\end{equation}
Its tangent contains material and geometric contributions,
\begin{equation}
 \bm K_M=\sum_g\Omega_g\left[
      \bm B_g^T\bm D_g\bm B_g+
      \sum_{\alpha=1}^{3}(s_g)_\alpha
       \frac{\partial^2(e_g)_\alpha}{\partial\bm U\,\partial\bm U}
       \right].
 \label{eq:macro_tangent}
\end{equation}
The first term uses the derivative of the deployed stress law; the second follows from nonlinear strain kinematics. The constitutive interface therefore requires both $\ssv_g$ and $\bm D_g$. It does not require the internal implementation to use a microscopic mesh, reduced coordinates, or a neural energy.

In FOM--FE$^2$, this interface invokes the full periodic solve. In the HPROM variants it invokes a reduced microscopic evaluator. In the learned constitutive variants it evaluates the stress and tangent map directly. The macroscopic mesh, loading and equilibrium algorithm are kept unchanged. The suffix FE$^2$ used in the archived model names identifies this common multiscale comparison; direct neural evaluations are not nested microscopic finite element solves.
